\documentclass[11pt]{amsart}
\usepackage[T1]{fontenc}
\usepackage{lmodern}
\usepackage{amsmath,amssymb,amsthm,mathtools}
\usepackage{xcolor}
\usepackage[colorlinks=true,linkcolor=blue!55!black,citecolor=blue!55!black,urlcolor=blue!55!black]{hyperref}

\newtheorem{theorem}{Theorem}[section]
\newtheorem{lemma}[theorem]{Lemma}
\newtheorem{corollary}[theorem]{Corollary}
\theoremstyle{remark}
\newtheorem{remark}[theorem]{Remark}

\newcommand{\LG}{\operatorname{LG}}
\newcommand{\Z}{\mathbb Z}

\title{Links--Gould positivity for two-block alternating
$2$-bridge knots}
\author{Zijian Zeng}
\address{Institute of Computer Science and Digital Innovation,
UCSI University, Kuala Lumpur 56000, Malaysia}
\email{zijianzeng@foxmail.com}
\email{1002266693@ucsiuniversity.edu.my}
\date{August 30, 2026}
\subjclass[2020]{57K14, 57K10}
\keywords{Links--Gould polynomial, alternating knot, $2$-bridge knot,
coefficient positivity, Ishii conjecture}
\hypersetup{
  pdftitle={Links--Gould positivity for two-block alternating 2-bridge knots},
  pdfauthor={Zijian Zeng}
}

\begin{document}

\begin{abstract}
Ishii conjectured that the coefficients of the Links--Gould polynomial of
an alternating knot have signs determined, up to one global sign, by total
degree parity.  We prove the positive normalization of this conjecture for
the two-block alternating $2$-bridge knots $D(p,q)$, with $p,q\geq1$, in
Kohli's notation.  More precisely, we give an explicit factorization of
$\LG(D(p,q);-x,-y)$ as a sum of products of Laurent polynomials with
nonnegative integer coefficients.  The proof is symbolic and uses Kohli's
three-term skein recurrence; exact computations serve only as independent
checks.  The result does not assert hole-freeness or log-concavity, and the
conjecture for arbitrary alternating knots remains open.
\end{abstract}

\maketitle

\section{Introduction}

Write the two-variable Links--Gould invariant in Ishii's normalization as
\[
  \LG(K;t_0,t_1)=\sum_{i,j}a_{ij}t_0^it_1^j.
\]
Ishii conjectured that, for every alternating knot $K$, coefficients whose
total degrees have the same parity have the same weak sign, while those of
opposite parity have opposite weak signs \cite{Ishii2006}.  Equivalently,
all nonzero coefficients of $\LG(K;-t_0,-t_1)$ have one global sign.

Ishii checked prime knots through ten crossings.  Harper, Kohli, Song, and
Tahar subsequently formulated stronger positivity and support conjectures,
reported verification of the weak sign statement through sixteen crossings,
and proved their conjectures for $(2,n)$-torus links and twist knots
\cite{HKST2025}.  Garoufalidis, Li, and Yu have since verified a still
stronger positivity, hole-free, and log-concavity statement for all
$51{,}280{,}976$ alternating knots with at most nineteen crossings
\cite{GLY2026}.  These are finite certifications and do not supply a proof
for arbitrary alternating knots.

Our result treats the two-block part of Kohli's alternating $2$-bridge
family.  Let $D(b_1,\ldots,b_m)$ denote the oriented $2$-bridge diagram of
\cite[Figures~3 and~4]{Kohli2018}.  When all $b_i$ are positive the diagram
is alternating, and it represents a knot precisely when $m$ is even
\cite[Remark~2.13]{Kohli2018}.  Thus $D(p,q)$ is an alternating knot for
all $p,q\geq1$.  The boundary subfamilies meet the previously studied twist
knots, while allowing both parameters to vary gives the full two-block
family.

\begin{theorem}\label{thm:main}
For every $p,q\geq1$,
\[
  \LG(D(p,q);-x,-y)\in\Z_{\geq0}[x^{\pm1},y^{\pm1}].
\]
Consequently the coefficients of $\LG(D(p,q);t_0,t_1)$ satisfy Ishii's
alternating-sign conjecture.
\end{theorem}

The proof converts the apparent rational coefficients in Kohli's recurrence
into positive finite sums.  No specialization, limit, or numerical sign
decision is involved.

\section{The two-block recurrence}

We use $\LG(\bigcirc)=1$.  Define
\begin{align*}
 a_1(n)&=\frac{(t_0t_1)^n-1}{t_0t_1-1},\\
 a_2(n)&=\frac{2n(t_0-1)(t_1-1)}{t_0t_1-1}
 -a_1(n)\left(
 \frac{(t_0t_1+1)(t_0-1)(t_1-1)}{t_0t_1-1}+1\right),\\
 a_3(n)&=(t_0-1)(t_1-1)a_1(n)+1.
\end{align*}
For a Laurent polynomial $f$, set
$\widetilde f(t_0,t_1)=f(t_0^{-1},t_1^{-1})$.  Corollary~2.12 and the proof
of Theorem~2.15 in \cite{Kohli2018} give
\begin{align}
 \LG(D(p))&=\widetilde a_1(p)+\widetilde a_2(p),
 \label{eq:oneblock}\\
 \LG(D(p,q))&=a_1(q)\LG(D(p-1))
 +a_2(q)\LG(D(p))+a_3(q).                 \label{eq:twoblock}
\end{align}
Here $D(0)$ is the split diagram and $\LG(D(0))=0$; the last term of
\eqref{eq:twoblock} uses the unknot normalization.  Although the displayed
definitions contain quotients, each $a_i(n)$ is a Laurent polynomial.

\section{Positive finite sums}

Set
\[
 U=xy,\qquad H^+=(1+x)(1+y),\qquad
 H^-=(1+x^{-1})(1+y^{-1}).
\]
For $n\geq0$, with an empty sum interpreted as zero, put
\[
 A_n^\pm=\sum_{k=0}^{n-1}U^{\pm k},\qquad
 E_n^\pm=\sum_{k=0}^{n-1}(2(n-k)-1)U^{\pm k}.
\]

\begin{lemma}\label{lem:sums}
For $n\geq1$, substitution $(t_0,t_1)=(-x,-y)$ gives
\begin{align}
 a_1(n)&=A_n^+,&
 -a_2(n)&=A_n^++H^+E_n^+,&
 a_3(n)&=1+H^+A_n^+ .                    \label{eq:ai-positive}
\end{align}
Moreover,
\begin{equation}
 E_n^--E_{n-1}^-=U^{-(n-1)}+2A_{n-1}^- .
 \label{eq:E-difference}
\end{equation}
\end{lemma}

\begin{proof}
The formulas for $a_1$ and $a_3$ follow immediately.  For $a_2$, the only
needed quotient identity is
\begin{equation}
 \frac{(U+1)A_n^+-2n}{U-1}=E_n^+.
 \label{eq:quotient}
\end{equation}
Indeed, multiplication by $U-1$ turns the right side into
\[
 U^n+2\sum_{k=1}^{n-1}U^k-(2n-1)
 =(U+1)\sum_{k=0}^{n-1}U^k-2n,
\]
which proves \eqref{eq:quotient} as a polynomial identity.  Substitution in
the definition of $a_2$ proves \eqref{eq:ai-positive}.  Finally,
\eqref{eq:E-difference} follows directly by subtracting the two finite sums.
\end{proof}

In \eqref{eq:ai-positive}, the arguments $(-x,-y)$ have been suppressed.
Inverting $x$ and $y$ gives the corresponding formulas with superscript
minus.

\section{Proof of the main theorem}

Let
\[
 G_p=-\LG(D(p);-x,-y).
\]
Equation \eqref{eq:oneblock} and Lemma~\ref{lem:sums} yield
\begin{equation}
 G_p=H^-E_p^- .                            \label{eq:one-positive}
\end{equation}
This also matches the boundary $p=0$, because both sides vanish.

Substitute \eqref{eq:ai-positive} and \eqref{eq:one-positive} into
\eqref{eq:twoblock}.  Before simplification the result is
\[
 -A_q^+H^-E_{p-1}^-
 +(A_q^++H^+E_q^+)H^-E_p^-
 +1+H^+A_q^+.
\]
Using \eqref{eq:E-difference} gives the exact factorization
\begin{equation}
\boxed{\begin{aligned}
\LG(D(p,q);-x,-y)={}&
A_q^+H^-\bigl(U^{-(p-1)}+2A_{p-1}^-\bigr)\\
&+H^+E_q^+H^-E_p^-+1+H^+A_q^+.
\end{aligned}}                             \label{eq:factorization}
\end{equation}
Every factor and every summand on the right belongs to
$\Z_{\geq0}[x^{\pm1},y^{\pm1}]$.  This proves the first assertion of
Theorem~\ref{thm:main}, including $p=1$ and $q=1$ by the empty-sum
convention.

Finally, if
$\LG(D(p,q);t_0,t_1)=\sum a_{ij}t_0^it_1^j$, coefficientwise positivity
after substituting $(-x,-y)$ says $(-1)^{i+j}a_{ij}\geq0$ for every
$i,j$.  Coefficients of equal total-degree parity therefore have
nonnegative products and coefficients of opposite parity have nonpositive
products, exactly as conjectured by Ishii.  \qed

\begin{remark}
For $p=q=1$, formula \eqref{eq:factorization} gives the figure-eight knot
and expands to
\begin{align*}
&2xy+3x+x/y+3y+7+3/y\\
&\hspace{35mm}+y/x+3/x+2/(xy),
\end{align*}
providing a convenient normalization check.
\end{remark}

\section{Exact verification and limitations}

The proof above is uniform in $p$ and $q$ and does not depend on computer
experiments.  For error detection, the ancillary archive contains two
independently written exact checks.  The first constructs
\eqref{eq:factorization} using integer dictionaries indexed by Laurent
exponent pairs.  The second starts from the rational definitions of
$a_1,a_2,a_3$, performs symbolic cancellation, reconstructs
\eqref{eq:oneblock}--\eqref{eq:twoblock}, and compares the result with
\eqref{eq:factorization}.  Two further implementations of the published
$R$-matrix, one in a quadratic Laurent quotient and one using the
radical-free $V_1$ matrix of \cite{Skein2026}, agree on the figure-eight
case.  All arithmetic in these checks is exact.

The result proves only the coefficient-sign part of the current conjectural
picture.  It does not prove that the support has no holes, nor either form of
log-concavity in \cite{HKST2025,GLY2026}.  It also does not settle Ishii's
conjecture for alternating knots outside the two-block $2$-bridge family.

\section*{Data and code availability}

The accompanying source archive includes the proof record and all exact
checkers.  From the archive root, the principal checks are run with
\begin{verbatim}
python3 certificate/double_twist_certificate.py 5 7
python3 certificate/verify_double_twist_sympy.py 5 7 \
  --max-identity-n 20
python3 code/lg_exact.py --self-test
python3 code/lg_v1_sympy.py --self-test
\end{verbatim}
The finite range in the symbolic regression test checks the implementation;
the all-parameter proof is Lemma~\ref{lem:sums} and
\eqref{eq:factorization}.

\bibliographystyle{amsplain}
\bibliography{references}

\end{document}
